\begin{abstract}
Multi-hop reasoning over \vic{real-life} knowledge graphs (KGs)  is a \vic{highly} challenging \vic{problem as traditional subgraph matching methods are not capable to deal with noise and missing information.} %because KGs are inevitably incomplete and noisy. 
\vicf{To address this problem, it has been recently  introduced a promising approach} based on \vic{jointly embedding logical queries and KGs} 
%embed entities and first-order logical (FOL) queries 
into a  low-dimensional space  to identify answer entities. 
However, \vic{existing proposals} ignore critical \vic{semantic} knowledge \vic{inherently available in KGs}, such as type information. %inside the KGs. %Besides, they have weak inductive reasoning ability.
%In order 
\vicf{To leverage  type information, %in the KGs, 
we propose a novel  \textbf{T}yp\textbf{E}-aware \textbf{M}essage \textbf{P}assing (\textsc{Temp}) model}, which enhances the entity and relation \jeff{representations} in queries, and simultaneously improves  generalization, deductive and inductive reasoning. %by leveraging type information and 
\vic{Remarkably,} \textsc{Temp} is a \vic{plug-and-play} model that can be easily incorporated into  existing \vic{embedding-based} models to improve their performance. %It has two distinct %effective
%type-driven mechanisms for query answering: i) TER: independently uses each neighbor type of an entity to enhance its representation; ii) TRR: match the entity content, relation content, and the aggregated type information of specific relation to further enhance the entity and relation representation. 
Extensive experiments on three real-world datasets %(FB15k, FB15K-237, and NELL995) 
demonstrate \vic{\textsc{Temp}'s} effectiveness. %\footnote{
 %Data, code, and an extended version with appendix  at \url{https://github.com/SXUNLP/QE-TEMP}
%}
\end{abstract}
